%% =====================================================================
%%  Unit 6 — A day in my life  (Raconter sa journée)
%%  Pistes : 6.01–6.04 dialogues · 6.05–6.06 vocabulaire · 6.07–6.08 grammaire
%%           6.09 prononciation · 6.10–6.13 exercices (textes en annexe)
%% =====================================================================

\newcommand{\otIllusUnVI}{\otIllusDuo{Monday/104/86/15/0.10,7:30/190/92/20/0.09,coffee/122/14/13/0.08,
  subway/150/8/11/0.08,work/196/22/12/0.07,sleep/100/50/12/0.07}{jack}{I get up at six!}{lea}{Six? Wow!}}

\unite[A day\\in my life]{6}{A day in my life}{Raconter sa journée}{otIllusUnVI}

\situation{Épisode 6}{Mardi, 9 h 15 au studio. Jack est arrivé depuis longtemps, Léa vient de
s'installer… et Emma est encore en retard. L'occasion de comparer les journées de chacun. Le soir,
Mrs.\ Lopez raconte à Léa la journée bien remplie de Sam.}

\begin{objectifs}
\item demander et dire l'heure ;
\item décrire votre routine, du réveil au coucher ;
\item utiliser le présent simple pour parler des habitudes ;
\item ajouter le \en{-s} de la 3\textsuperscript{e} personne : \en{he works}, \en{she goes} ;
\item nommer les jours de la semaine et les moments de la journée.
\end{objectifs}

\auprogramme
  {Décrire sa routine, dire l'heure}
  {Le présent simple (affirmatif)\par\vspace{0.6mm}Le \en{-s} de la 3\textsuperscript{e} personne}
  {Activités quotidiennes, l'heure, les jours : 38 expressions}
  {Le \en{-s} de la 3\textsuperscript{e} personne}
  {\en{a.m.} et \en{p.m.}, la journée de travail}

\semaine{casque/otNavy/Dialogues/Écouter et lire les deux scènes,
  bulle/otGreenDark/Vocabulaire/Apprendre les 38 expressions,
  blocs/otViolet/Grammaire/Présent simple et heure,
  sons/otOrangeDark/Sons et culture/Le -s et la journée américaine,
  crayon/otNavy/Exercices/12 exercices corrigés,
  coche/otGreenDark/Bilan/L'essentiel et le test express}

\clearpage
%% ---------------------------------------------------------------------
\rubrique{ecouter}{Matin et soir}

\begin{dialogue}{1}{I get up at six}{Mardi, 9 h 15, au studio. Léa et Jack sont à leur bureau.
  Emma n'est pas encore là.}{6.01}{6.02}
\dit{lea}{Jack, \hl{what time is it}?\note{1}}{Jack, quelle heure est-il ?}
\dit{jack}{\hl{It's nine fifteen}. Emma is late again!}{Il est 9 h 15. Emma est encore en retard !}
\dit{lea}{Ha ha! And you're always early!}{Ha ha ! Et toi, tu es toujours en avance !}
\dit{jack}{Yes. \hl{I get up} at six every day.\note{2}}{Oui. Je me lève à six heures tous les jours.}
\dit{lea}{At six? Wow!}{À six heures ? Waouh !}
\dit{jack}{\hl{I go} to the gym, then \hl{I have breakfast} at home. \hl{I take} the subway at eight thirty.}{Je vais à la salle de sport, puis je prends mon petit-déjeuner chez moi. Je prends le métro à 8 h 30.}
\dit{lea}{Not me! I get up at seven thirty, \hl{I drink} a coffee and \hl{I walk} to the studio.}{Pas moi ! Je me lève à 7 h 30, je bois un café et je vais au studio à pied.}
\dit{emma}{Good morning! Sorry, I'm late!}{Bonjour ! Désolée, je suis en retard !}
\dit{jack}{Ha! \hl{Emma gets up} at nine o'clock!\note{3}}{Ha ! Emma, elle, se lève à neuf heures !}
\dit{emma}{That's not true! I get up at eight. But my train is always late!}{C'est faux ! Je me lève à huit heures. Mais mon métro est toujours en retard !}
\end{dialogue}

\begin{comprendre}
\pointnote{1} \en{\textbf{What time is it?}} « Quelle heure est-il ? ». On répond avec \en{It's…} :
  \en{It's nine fifteen} (9 h 15), \en{It's nine o'clock} (9 h pile).
\pointnote{2} \en{\textbf{I get up at six.}} Pour une habitude, l'anglais emploie le
  \textbf{présent simple} : le verbe à sa forme de base, sans \en{to}. \en{every day} = « tous les
  jours » ; \en{at six} = « à six heures ».
\pointnote{3} \en{\textbf{Emma gets up…}} À la 3\textsuperscript{e} personne du singulier
  (\en{he}, \en{she}, \en{it}, ou un nom comme \en{Emma}), le verbe prend un \en{-s}. C'est la seule
  terminaison du présent simple !
\end{comprendre}

\begin{dialogue}{2}{He has a long day}{Mardi soir, dans l'escalier de l'immeuble. Mrs.\ Lopez parle
  de Sam à Léa.}{6.03}{6.04}
\dit{lopez}{Sam is always so busy!}{Sam est toujours tellement occupé !}
\dit{lea}{Really?}{Ah bon ?}
\dit{lopez}{Yes. \hl{He works} in a coffee shop in the morning. \hl{He starts} at six.}{Oui. Il travaille dans un café le matin. Il commence à six heures.}
\dit{lea}{At six? That's early!}{À six heures ? C'est tôt !}
\dit{lopez}{Then \hl{he goes} to his music school in the afternoon.\note{4}}{Ensuite, il va à son école de musique l'après-midi.}
\dit{lea}{And in the evening?}{Et le soir ?}
\dit{lopez}{\hl{He plays} the guitar with his band. \hl{He finishes} at midnight!}{Il joue de la guitare avec son groupe. Il finit à minuit !}
\dit{lea}{Wow, \hl{he has} a long day!\note{5}}{Eh bien, il a une longue journée !}
\dit{lopez}{Yes, and on Sundays, \hl{he sleeps} all day!}{Oui, et le dimanche, il dort toute la journée !}
\end{dialogue}

\begin{comprendre}
\pointnote{4} \en{\textbf{he goes}, \textbf{he finishes}} Après \en{o}, \en{s}, \en{sh}, \en{ch} ou
  \en{x}, la 3\textsuperscript{e} personne s'écrit \en{-es} : \en{go} → \en{goes}, \en{finish} →
  \en{finishes}, \en{watch} → \en{watches}.
\pointnote{5} \en{\textbf{he has}} \en{have} est irrégulier : \en{I have}, mais \en{he has}.
  \en{on Sundays} = « le dimanche » (tous les dimanches) ; \en{all day} = « toute la journée ».
\end{comprendre}

%% ---------------------------------------------------------------------
\rubrique{mots}{Les mots de l'unité}

Voici les 38 expressions de l'unité. Les verbes sont donnés sans \en{to} : c'est la forme qu'on
emploie au présent simple, avec \en{I}, \en{you}, \en{we} et \en{they}.

\begin{vocabulaire}{Ma journée}{6.05}
\categorie{Le matin}
\mot{wake up}{ˌweɪk ˈʌp}{se réveiller}
\mot{get up}{ˌɡet ˈʌp}{se lever}
\mot{take a shower}{ˌteɪk ə ˈʃaʊɚ}{prendre une douche}
\mot{have breakfast}{ˌhæv ˈbrekfəst}{prendre le petit-déjeuner}
\mot{take the subway}{ˌteɪk ðə ˈsʌbweɪ}{prendre le métro}
\mot{walk}{wɑːk}{marcher, aller à pied}
\mot{go to work}{ˌɡoʊ tə ˈwɝːk}{aller au travail}
\mot{start work}{ˌstɑːrt ˈwɝːk}{commencer le travail}
\categorie{L'après-midi et le soir}
\mot{have lunch}{ˌhæv ˈlʌntʃ}{déjeuner}
\mot{finish work}{ˌfɪnɪʃ ˈwɝːk}{finir le travail}
\mot{go home}{ˌɡoʊ ˈhoʊm}{rentrer à la maison}
\mot{cook dinner}{ˌkʊk ˈdɪnɚ}{préparer le dîner}
\mot{watch TV}{ˌwɑːtʃ ˌtiːˈviː}{regarder la télé}
\mot{play the guitar}{ˌpleɪ ðə ɡɪˈtɑːr}{jouer de la guitare}
\mot{go to bed}{ˌɡoʊ tə ˈbed}{aller se coucher}
\mot{sleep}{sliːp}{dormir}
\categorie{Autres verbes utiles}
\mot{drink a coffee}{ˌdrɪŋk ə ˈkɑːfi}{boire un café}
\mot{study}{ˈstʌdi}{étudier}
\end{vocabulaire}

\begin{vocabulaire}{L'heure et les jours}{6.06}
\categorie{L'heure}
\mot{What time is it?}{ˌwʌt ˈtaɪm ɪz ɪt}{Quelle heure est-il ?}
\mot{It's nine o'clock.}{ɪts ˌnaɪn əˈklɑːk}{Il est neuf heures.}
\mot{It's nine fifteen.}{ɪts ˌnaɪn fɪfˈtiːn}{Il est neuf heures et quart.}
\mot{It's nine thirty.}{ɪts ˌnaɪn ˈθɝːti}{Il est neuf heures et demie.}
\mot{at six}{æt ˈsɪks}{à six heures}
\mot{noon · midnight}{nuːn · ˈmɪdnaɪt}{midi · minuit}
\categorie{Les moments}
\mot{in the morning}{ɪn ðə ˈmɔːrnɪŋ}{le matin}
\mot{in the afternoon}{ɪn ði ˌæftɚˈnuːn}{l'après-midi}
\mot{in the evening}{ɪn ði ˈiːvnɪŋ}{le soir}
\mot{every day}{ˌevri ˈdeɪ}{tous les jours}
\mot{early · late}{ˈɝːli · leɪt}{tôt, en avance · tard, en retard}
\mot{always}{ˈɑːlweɪz}{toujours}
\categorie{Les jours de la semaine}
\mot{Monday}{ˈmʌndeɪ}{lundi}
\mot{Tuesday}{ˈtuːzdeɪ}{mardi}
\mot{Wednesday}{ˈwenzdeɪ}{mercredi}
\mot{Thursday}{ˈθɝːzdeɪ}{jeudi}
\mot{Friday}{ˈfraɪdeɪ}{vendredi}
\mot{Saturday}{ˈsætɚdeɪ}{samedi}
\mot{Sunday}{ˈsʌndeɪ}{dimanche}
\mot{on Sundays}{ɑːn ˈsʌndeɪz}{le dimanche (tous les dimanches)}
\end{vocabulaire}

\begin{astuce}[Les jours en anglais]
Ils prennent toujours une \textbf{majuscule} : \en{Monday}, \en{Sunday}. Pour dire « le lundi »
(chaque lundi), on emploie \en{on} + le jour au pluriel : \en{on Mondays}. Et \en{Wednesday} se
prononce \phon{ˈwenzdeɪ} : le premier \en{d} ne s'entend pas.
\end{astuce}

%% ---------------------------------------------------------------------
\rubrique{grammaire}{Le présent simple et l'heure}

\otIntertitre[otViolet]{1 · Le présent simple : la forme affirmative}
Le présent simple sert à parler des \textbf{habitudes}, des \textbf{routines} et des \textbf{faits
généraux}. Il se forme avec la base verbale, et une seule terminaison : le \en{-s} de la
3\textsuperscript{e} personne du singulier.

\begin{center}\otTableau
\begin{tabular}{@{}>{\enOn}l<{\enOff}@{\hspace{10mm}}>{\enOn}l<{\enOff}@{\hspace{10mm}}l@{}}
\otTH{Sujet} & \otTH{Verbe} & \otTH{Français}\\\midrule
I, you, we, they        & work          & je travaille, tu travailles…\\
he, she, it, Léa, Jack  & work\textbf{s} & il / elle travaille\\\bottomrule
\end{tabular}\end{center}

\begin{paire}
\begin{regle}[L'orthographe du \en{-s}]
\exen{work → works}{cas général}
\exen{go → goes, watch → watches}{après \en{o}, \en{s}, \en{sh}, \en{ch}, \en{x} : \en{-es}}
\exen{study → studies}{consonne + \en{y} : \en{-ies}}
\exen{play → plays}{voyelle + \en{y} : \en{-s}}
\exen{have → has}{irrégulier}
\end{regle}
\begin{regle}[Les mots de l'habitude]
\exen{every day}{tous les jours}
\exen{in the morning}{le matin}
\exen{on Mondays}{le lundi}
\exen{at seven}{à sept heures}
\exen{always}{toujours (avant le verbe)}
\end{regle}
\end{paire}

\blocaudio{6.07}{Écoutez et répétez}{Une phrase avec \en{I}, puis la même avec \en{he} ou \en{she} :
écoutez bien le \en{-s} de la 3\textsuperscript{e} personne.}

\otIntertitre[otViolet]{2 · Dire l'heure}
\begin{center}
\begin{tikzpicture}
  \foreach \h/\m/\x/\lab/\uk in {9/0/0/nine o'clock/nine o'clock,9/15/38/nine fifteen/a quarter past nine,
      9/30/76/nine thirty/half past nine,9/45/114/nine forty-five/a quarter to ten}{%
    \otDessinHorloge{(\x mm,0)}{9mm}{\h}{\m}{otViolet}
    \node[anchor=north,font=\otSemi\fontsize{9.4}{11}\selectfont] at (\x mm,-11mm) {\enb{\lab}};
    \node[anchor=north,font=\otRegular\fontsize{8}{10}\selectfont,text=otMuted] at (\x mm,-16mm) {\enb{\uk}};}
\end{tikzpicture}
\end{center}
{\fontsize{10}{13}\selectfont La première façon, chiffre par chiffre, est la plus simple et la plus
courante aux États-Unis. La seconde (\en{past} = après, \en{to} = avant, \en{a quarter} = un quart,
\en{half} = demi) s'entend partout, surtout au Royaume-Uni. Pour dire « à » une heure : \en{at} :
\en{at nine thirty}.\par}

\blocaudio{6.08}{Écoutez et répétez}{Six heures, dites des deux façons quand c'est possible.}

\begin{attention}[Les pièges des francophones]
\fauxjuste{He work in a café.}{He works in a café.}
\fauxjuste{She gos to work.}{She goes to work.}
\fauxjuste{He haves breakfast.}{He has breakfast.}
\fauxjuste{I am get up at six.}{I get up at six.}
\fauxjuste{It's 9 hours.}{It's nine o'clock.}
\par\vspace{1mm}{\fontsize{10}{12.6}\selectfont Au présent simple, pas de \en{be} devant le verbe :
« je me lève » se dit \en{I get up}, jamais \en{I am get up}.\par}
\end{attention}

%% ---------------------------------------------------------------------
\rubrique[Le -s de la 3e personne]{sons}{Le \en{-s} de la 3\textsuperscript{e} personne}

Bonne nouvelle : le \en{-s} de \en{he works} se prononce exactement comme celui du pluriel
(unité 5) : \phon{s}, \phon{z} ou \phon{ɪz}. Trois verbes très courants font exception :
\en{does} \phon{dʌz}, \en{says} \phon{sez} et \en{has} \phon{hæz}.

\begin{paires}{6.09}{Écoutez le verbe avec \en{I}, puis avec \en{he}, et répétez. Ne laissez jamais
tomber le \en{-s} final : c'est lui qui montre qu'on parle d'une seule personne.}
\sonpaire{work}{wɝːk}{\phon{s}}{work\otH{s}}{wɝːks}{il travaille}
\sonpaire{sleep}{sliːp}{\phon{s}}{sleep\otH{s}}{sliːps}{il dort}
\sonpaire{play}{pleɪ}{\phon{z}}{play\otH{s}}{pleɪz}{il joue}
\sonpaire{go}{ɡoʊ}{\phon{z}}{go\otH{es}}{ɡoʊz}{il va}
\sonpaire{watch}{wɑːtʃ}{\phon{ɪz}}{watch\otH{es}}{ˈwɑːtʃɪz}{il regarde}
\sonpaire{finish}{ˈfɪnɪʃ}{\phon{ɪz}}{finish\otH{es}}{ˈfɪnɪʃɪz}{il finit}
\sonpaire{do}{duː}{\phon{uː}}{do\otH{es}}{dʌz}{il fait}
\sonpaire{say}{seɪ}{\phon{eɪ}}{say\otH{s}}{sez}{il dit}
\end{paires}

%% ---------------------------------------------------------------------
\rubrique{culture}{La journée à l'américaine}

\begin{paire}
\begin{culture}[\en{a.m.} et \en{p.m.}]
On ne compte pas les heures jusqu'à 24 : on dit \en{seven a.m.} (7 h du matin) et \en{seven p.m.}
(19 h). \en{a.m.} vient du latin \en{ante meridiem}, « avant midi ». L'heure « militaire » (19:00)
ne s'emploie que dans l'armée, les hôpitaux… et les horaires de train.
\end{culture}
\begin{culture}[\en{A nine-to-five job}]
La journée de bureau type va de 9 heures à 17 heures. La pause déjeuner est courte : beaucoup
d'Américains déjeunent d'un sandwich, à leur bureau (\en{lunch at my desk}), et dînent tôt, vers
18 h 30.
\end{culture}
\end{paire}

\begin{culture}[La semaine commence le dimanche]
Sur les calendriers américains, la première colonne est celle du dimanche ! Le week-end
(\en{the weekend}) désigne le samedi et le dimanche, et l'on souhaite \en{Have a nice weekend!}
le vendredi en partant.
\end{culture}

%% ---------------------------------------------------------------------
\rubrique{pratique}{Exercices}

\exercice[lire,ecrire,niveau=1]{Quelle heure est-il ?}{Écrivez l'heure en anglais, chiffre par
chiffre : \en{It's seven o'clock.}, \en{It's eight fifteen.}…}
\begin{tcbraster}[raster columns=3,raster equal height=rows,raster column skip=3mm,raster row skip=3mm,raster before skip=1mm]
\otCarteHeure{7}{0}{7 h 00}{}{soleil}{It's seven o'clock.}
\otCarteHeure{8}{15}{8 h 15}{}{soleil}{It's eight fifteen.}
\otCarteHeure{10}{30}{10 h 30}{}{soleil}{It's ten thirty.}
\otCarteHeure{12}{45}{12 h 45}{}{soleil}{It's twelve forty-five.}
\otCarteHeure{3}{20}{15 h 20}{}{soleil}{It's three twenty.}
\otCarteHeure{6}{50}{18 h 50}{}{lune}{It's six fifty.}
\end{tcbraster}

\exercice[ecrire,niveau=1]{À la 3\textsuperscript{e} personne}{Écrivez le verbe à la 3\textsuperscript{e}
personne du singulier.}
\begin{questions}[2]
\item I work → he \trou[26mm]{works}
\item I go → she \trou[26mm]{goes}
\item I watch → he \trou[26mm]{watches}
\item I study → she \trou[26mm]{studies}
\item I have → he \trou[26mm]{has}
\item I play → she \trou[26mm]{plays}
\item I finish → he \trou[26mm]{finishes}
\item I do → she \trou[26mm]{does}
\end{questions}

\exercice[lire,ecrire,niveau=2]{La journée de Sam}{Complétez avec le verbe entre parenthèses, au présent simple.}
\begin{anglais}
\begin{tcolorbox}[otBase,colback=otNavyTint,colframe=otNavyTint,fontupper=\fontsize{11}{19}\selectfont]
Sam \trou[22mm]{gets up} \textit{(get up)} at five. He \trou[22mm]{takes} \textit{(take)} a shower and
he \trou[22mm]{goes} \textit{(go)} to the coffee shop. He \trou[22mm]{starts} \textit{(start)} work at six.
In the afternoon, he \trou[22mm]{studies} \textit{(study)} music. In the evening, he
\trou[22mm]{plays} \textit{(play)} the guitar with his band. He \trou[22mm]{has} \textit{(have)} dinner
at eleven and he \trou[22mm]{finishes} \textit{(finish)} at midnight.
\end{tcolorbox}
\end{anglais}

\exercice[ecrire,niveau=1]{Le bon verbe}{Choisissez la bonne forme et écrivez-la.}
\begin{questions}
\item My mother \trou[22mm]{works} in a hospital. \textit{(work / works)}
\item Jack and Emma \trou[22mm]{take} the subway. \textit{(take / takes)}
\item I \trou[22mm]{drink} a coffee every morning. \textit{(drink / drinks)}
\item Léa \trou[22mm]{walks} to the studio. \textit{(walk / walks)}
\item We \trou[22mm]{have} lunch at one. \textit{(have / has)}
\item Mrs.\ Lopez \trou[22mm]{watches} TV in the evening. \textit{(watch / watches)}
\end{questions}

\exercice[ecrire,niveau=1]{Les jours de la semaine}{Complétez.}
\begin{questions}
\item Monday, \trou[30mm]{Tuesday}, Wednesday, \trou[30mm]{Thursday}, Friday
\item The day after Friday is \trou[30mm]{Saturday}.
\item The day before Monday is \trou[30mm]{Sunday}.
\item « mercredi » = \trou[30mm]{Wednesday}
\item « le lundi » (tous les lundis) = \trou[30mm]{on Mondays}
\end{questions}

\exercice[lire,niveau=2]{La journée de Léa}{Remettez les phrases dans l'ordre de la journée : écrivez
les lettres.}
\begin{anglais}
\begin{questions}[2][english]
\item[a.] She has lunch at one o'clock.
\item[b.] She gets up at seven thirty.
\item[c.] She goes to bed at eleven.
\item[d.] She starts work at nine.
\item[e.] She has breakfast at eight.
\item[f.] She finishes work at six.
\end{questions}
\end{anglais}
\par\vspace{2mm}\noindent{\otSemi\fontsize{10}{12}\selectfont Ordre :}\hspace{4mm}%
\trou[9mm]{b}\hspace{3mm}\trou[9mm]{e}\hspace{3mm}\trou[9mm]{d}\hspace{3mm}\trou[9mm]{a}\hspace{3mm}\trou[9mm]{f}\hspace{3mm}\trou[9mm]{c}

\exercice[ecouter,niveau=1,piste=6.10]{\phon{s}, \phon{z} ou \phon{ɪz} ?}{Écoutez chaque verbe, dit deux fois.
Cochez le son de la terminaison.}
\begin{questions}[3]
\item works \hfill \choix{\phon{s}} \choix{\phon{z}} \choix{\phon{ɪz}}
\item plays \hfill \choix{\phon{s}} \choix{\phon{z}} \choix{\phon{ɪz}}
\item watches \hfill \choix{\phon{s}} \choix{\phon{z}} \choix{\phon{ɪz}}
\item sleeps \hfill \choix{\phon{s}} \choix{\phon{z}} \choix{\phon{ɪz}}
\item goes \hfill \choix{\phon{s}} \choix{\phon{z}} \choix{\phon{ɪz}}
\item finishes \hfill \choix{\phon{s}} \choix{\phon{z}} \choix{\phon{ɪz}}
\item eats \hfill \choix{\phon{s}} \choix{\phon{z}} \choix{\phon{ɪz}}
\item reads \hfill \choix{\phon{s}} \choix{\phon{z}} \choix{\phon{ɪz}}
\item dances \hfill \choix{\phon{s}} \choix{\phon{z}} \choix{\phon{ɪz}}
\end{questions}

\exercice[ecouter,ecrire,niveau=2,piste=6.11]{La journée d'Emma}{Emma raconte sa journée. Écoutez et
notez l'heure de chaque activité.}
\begin{center}\otTableau
\begin{tabular}{@{}l@{\hspace{14mm}}l@{}}
\otTH{Activité} & \otTH{Heure}\\\midrule
elle se lève                 & \trou[22mm]{8:00}\\[0.6mm]
elle prend le métro          & \trou[22mm]{8:40}\\[0.6mm]
elle commence le travail     & \trou[22mm]{9:15}\\[0.6mm]
elle déjeune                 & \trou[22mm]{1:30}\\[0.6mm]
elle finit le travail        & \trou[22mm]{6:00}\\[0.6mm]
elle va au yoga (le mardi)   & \trou[22mm]{7:00}\\[0.6mm]
elle se couche               & \trou[22mm]{11:30}\\\bottomrule
\end{tabular}\end{center}

\exercice[ecouter,ecrire,niveau=2,piste=6.12]{Dictée à trous}{Léa raconte à Sam la journée de Jack.
Écoutez et complétez.}
\begin{anglais}
\begin{tcolorbox}[otBase,colback=white,colframe=otRule,boxrule=0.5pt,fontupper=\fontsize{11}{19}\selectfont]
Jack \trou{gets} up at six. He \trou{goes} to the gym, then he \trou{has} breakfast at home.
He \trou{takes} the subway at eight thirty and he \trou{starts} work at nine. He \trou{has}
lunch at his desk. In the evening, he \trou{watches} basketball on TV. He \trou{goes} to bed at ten!
\end{tcolorbox}
\end{anglais}

\exercice[ecrire,niveau=2]{Traduisez}{Traduisez en anglais.}
\begin{questions}[1][french]
\item Quelle heure est-il ? — Il est huit heures et demie.\lignes{1}
\item Je me lève à sept heures tous les jours.\lignes{1}
\item Emma prend le métro.\lignes{1}
\item Il finit le travail à six heures.\lignes{1}
\item Le samedi, je dors toute la journée.\lignes{1}
\item Ma sœur étudie la musique.\lignes{1}
\item Nous déjeunons à midi.\lignes{1}
\item Jack regarde la télé le soir.\lignes{1}
\end{questions}

\exercice[ecrire,niveau=3]{Votre journée}{Racontez une journée habituelle en six phrases, avec des heures :
\en{I get up at…}, \en{I have breakfast…}, \en{In the evening, I…}}
\lignes{6}

\exercice[ecouter,parler,niveau=3,piste=6.13]{À vous de parler !}{Jack s'intéresse à votre
journée. Répondez-lui \textbf{à voix haute} pendant chaque pause, en deux ou trois phrases.}
\begin{center}\otTableau
\begin{tabular}{@{}>{\otSemi\fontsize{8.6}{10}\selectfont\color{otNavy}}l@{\hspace{3mm}}>{\enOn}p{70mm}<{\enOff}@{\hspace{6mm}}p{78mm}@{}}
 & \otTH{Jack dit\ldots} & \otTH{Vous\ldots}\\\midrule
1. & Hi! What time is it? & donnez l'heure (celle de votre montre, ou 9 h 15).\\
2. & Tell me about your mornings. & dites à quelle heure vous vous levez et ce que vous faites.\\
3. & And in the afternoon? & décrivez vos après-midi.\\
4. & And in the evening? & dites ce que vous faites le soir.\\
5. & Wow, that's a busy day! & réagissez simplement.\\\bottomrule
\end{tabular}\end{center}

%% ---------------------------------------------------------------------
\rubrique{bilan}{Faites le point}

\begin{essentiel}
\begin{minipage}[t]{0.52\linewidth}
\otEssTitre{Les phrases clés}
\renewcommand{\arraystretch}{1.32}
\begin{tabular}{@{}>{\enOn\otSemi\fontsize{9.4}{11}\selectfont}l<{\enOff}@{\hspace{3mm}}>{\fontsize{9.2}{11}\selectfont\color{white!78!otNavy}}l@{}}
What time is it?            & Quelle heure est-il ?\\
It's nine thirty.           & Il est 9 h 30.\\
I get up at seven.          & Je me lève à 7 h.\\
She takes the subway.       & Elle prend le métro.\\
He has lunch at one.        & Il déjeune à 13 h.\\
on Mondays · every day      & le lundi · tous les jours\\
\end{tabular}
\end{minipage}\hfill
\begin{minipage}[t]{0.42\linewidth}
\otEssTitre{Présent simple}
{\enb{\fontsize{9.4}{12}\selectfont I / you / we / they \textbf{work}\par he / she / it \textbf{works}}\par}
\otEssTitre{Orthographe}
{\fontsize{9.2}{11.6}\selectfont \en{goes}, \en{watches}, \en{studies}, \en{has}\par}
\otEssTitre{L'heure}
{\fontsize{9.2}{11.6}\selectfont \en{It's nine fifteen} ; \en{at} + heure ; \en{a.m.} / \en{p.m.}\par}
\otEssTitre{Prononciation}
{\fontsize{9.2}{11.6}\selectfont \en{works} \phon{s} · \en{plays} \phon{z} · \en{watches} \phon{ɪz} · \en{does} \phon{dʌz}\par}
\end{minipage}
\end{essentiel}

\begin{jesais}
\item demander et dire l'heure, et dire « à » une heure avec \en{at} ;
\item décrire ma routine, du matin au soir ;
\item conjuguer un verbe au présent simple, avec le \en{-s} de la 3\textsuperscript{e} personne ;
\item écrire correctement \en{goes}, \en{watches}, \en{studies}, \en{has} ;
\item nommer les jours de la semaine et dire \en{on Mondays} ;
\item prononcer \en{works}, \en{plays}, \en{watches} et \en{does}.
\end{jesais}

\clearpage
\begin{testexpress}
\qcm[b]{« Quelle heure est-il ? » se dit :}{\en{What hour is it?}}{\en{What time is it?}}{\en{How time is it?}}
\qcm[a]{9 h 30 se dit :}{\en{nine thirty}}{\en{nine three}}{\en{thirty nine}}
\qcm[c]{Complétez : \en{She \rule[-0.6mm]{9mm}{0.5pt} to work.}}{\en{go}}{\en{gos}}{\en{goes}}
\qcm[a]{Complétez : \en{They \rule[-0.6mm]{9mm}{0.5pt} breakfast at eight.}}{\en{have}}{\en{has}}{\en{haves}}
\qcm[b]{« Je me lève à six heures » se dit :}{\en{I am get up at six.}}{\en{I get up at six.}}{\en{I gets up at six.}}
\qcm[c]{La 3\textsuperscript{e} personne de \en{study} est :}{\en{studys}}{\en{studyes}}{\en{studies}}
\qcm[b]{« le dimanche » (chaque dimanche) se dit :}{\en{the Sunday}}{\en{on Sundays}}{\en{in Sunday}}
\qcm[a]{\en{does} se prononce :}{\phon{dʌz}}{\phon{duːz}}{\phon{doʊz}}
\qcm[c]{« mercredi » se dit :}{\en{Wendsday}}{\en{Thursday}}{\en{Wednesday}}
\qcm[b]{\en{seven p.m.}, c'est :}{7 h du matin}{19 h}{minuit moins sept}
\end{testexpress}

\noindent{\fontsize{9.6}{12}\selectfont\color{otMuted}Niveau visé (CECRL, niveau A1) : « peut
décrire ses activités quotidiennes et donner l'heure ». Les corrigés de l'unité sont
page~\pageref{cor-u6-1} ; les textes des pistes d'exercices, page~\pageref{trans-6.10}.\par}
